\documentclass[a4paper,article,14pt]{extarticle}

\usepackage{../spbu_diploma/spbudiploma_tempora}

\usepackage{euscript}
\usepackage{longtable}
\usepackage{makecell}
\usepackage[pdftex]{graphicx}
\usepackage{amsthm,amssymb,amsmath}
\usepackage{textcomp}
\usepackage{array}

\begin{document}

\begin{titlepage}
\begin{center}
\textbf{Санкт--Петербургский государственный университет}

\vspace{35mm}

\textbf{\large Выпускная квалификационная работа}\\[3mm]
\textbf{\large Разработка рекомендательной системы севооборота}

\vspace{20mm}
Уровень образования: Магистратура\\
Направление подготовки: Программирование и информационные технологии\\
Образовательная программа: Технологии искусственного интеллекта и Big Data

\vspace{25mm}
\begin{flushright}
\begin{minipage}[t]{0.65\textwidth}
Автор: Корельский Дмитрий Сергеевич\\
Научный руководитель: Митрофанов Евгений Павлович
\end{minipage}
\end{flushright}

\vfill
Санкт-Петербург\\
\the\year{} г.
\end{center}
\end{titlepage}

\tableofcontents
\pagebreak

\specialsection{Введение}

Задача выбора следующей культуры на сельскохозяйственном поле относится к числу практических задач, в которых исторические данные могут быть полезны для поддержки принятия решений. На практике специалисту часто важно видеть не только один наиболее вероятный вариант, но и короткий список допустимых кандидатов, сопоставимый с историей поля, региональным контекстом и наблюдаемыми закономерностями переходов между культурами. Поэтому данная работа рассматривает задачу не как построение жесткого предписания, а как разработку объяснимого рекомендательного прототипа.

Тема работы --- рекомендательная система по севообороту на основе истории культур по полям. В рамках диплома система формирует ранжированный список вероятных следующих культур для конкретного поля. Такой результат полезен как shortlist для дальнейшего экспертного рассмотрения: пользователь видит несколько кандидатов, оценку уверенности модели и объяснительный контекст. Работа сознательно не позиционируется как полноценный агрономический оптимизатор, поскольку в ней не моделируются экономические ограничения, погодные прогнозы, почвенные факторы, доступность техники, хозяйственные планы и другие условия, необходимые для гарантированного выбора оптимального севооборота.

Объектом исследования является процесс построения рекомендаций следующей культуры по историческим данным о полях. Предметом исследования являются методы подготовки временных окон, построения табличной модели классификации, формирования Top-k рекомендаций и объяснения кандидатов через дополнительные аналитические слои.

Цель работы состоит в построении воспроизводимого исследовательского прототипа, который по истории культур на поле и дополнительным признакам поля или региона формирует ранжированный список вероятных следующих культур, сопровождаемый оценкой уверенности и интерпретацией.

Для достижения цели решаются следующие задачи:
\begin{enumerate}
    \item подготовить данные USDA NASS CSB/CDL и привести их к формату, пригодному для моделирования истории культур;
    \item сформировать оконную выборку вида \texttt{history\_1, history\_2, history\_3 -> target};
    \item зафиксировать воспроизводимое разделение train, validation и test с учетом группировки по \texttt{CSBID};
    \item обучить и оценить финальную baseline-модель CatBoost;
    \item перейти от Top-1 классификации к Top-k рекомендательному сценарию;
    \item оценить confidence layer и калибровку вероятностей;
    \item сравнить CatBoost с простым Markov-1 baseline;
    \item построить слой объяснимости на основе knowledge graph и правил;
    \item показать пространственную демонстрацию рекомендаций на карте.
\end{enumerate}

Практическая значимость работы состоит в том, что построен полный notebook-first и artifact-driven контур: от подготовки данных и фиксированного split до модели, Top-k рекомендаций, confidence analysis, knowledge graph explainability и пространственной демонстрации. Финальным предиктивным ядром диплома является \texttt{baseline} CatBoost. Ветка \texttt{improved} рассматривается только как экспериментальная и не заменяет финальную модель.

Работа имеет следующую структуру. В первой главе задается постановка задачи и границы применимости. Во второй главе описывается архитектура решения и проектные контракты. Третья глава посвящена данным и подготовке выборки. В четвертой главе рассматриваются baseline-подходы и финальная CatBoost-модель. В пятой и шестой главах описываются рекомендательный слой, Top-k оценка, confidence buckets и калибровка. Седьмая глава посвящена explainability через knowledge graph. Восьмая глава описывает case studies и пространственную демонстрацию. Девятая глава фиксирует ограничения и направления дальнейшего развития.

\section{Постановка задачи и позиционирование работы}

\subsection{Постановка задачи}

В работе решается задача рекомендации следующей культуры для поля на основе его недавней истории и контекстных признаков. Исходная информация по полю включает последовательность культур за несколько лет, региональные признаки, площадь и координаты. На выходе система выдает вероятности по классам следующей культуры и преобразует их в ранжированный список кандидатов.

Формально математическая основа задачи --- многоклассовая классификация. Для объекта $x$ модель оценивает распределение вероятностей по множеству классов культур $C$:
\begin{equation}
    p(c \mid x), \quad c \in C.
\end{equation}
Затем классы сортируются по убыванию вероятности, и пользователь получает Top-k список. В дипломе особое внимание уделяется Top-3 и Top-5 рекомендациям, поскольку именно они соответствуют прикладной логике shortlist.

Входные признаки финальной модели включают:
\begin{itemize}
    \ITEM \texttt{history\_1}, \texttt{history\_2}, \texttt{history\_3} --- трехлетнюю историю культур;
    \ITEM \texttt{STATEFIPS}, \texttt{CNTYFIPS} --- региональный контекст;
    \ITEM \texttt{CSBACRES} --- площадь поля;
    \ITEM \texttt{INSIDE\_X}, \texttt{INSIDE\_Y} --- пространственные координаты.
\end{itemize}

Выходом модели являются вероятности по девяти финальным классам: \texttt{corn}, \texttt{cotton}, \texttt{fallow}, \texttt{forage\_hay}, \texttt{legumes}, \texttt{other\_cereals}, \texttt{sorghum}, \texttt{soybeans}, \texttt{wheat}. Прикладным результатом является не только наиболее вероятный класс, но и Top-k список с confidence-сигналом и объяснением.

Важно зафиксировать границу постановки. Система не решает задачу оптимального севооборота в полном агрономическом смысле. Она не рассчитывает прибыльность, не учитывает доступность ресурсов, не моделирует долгосрочный план хозяйства и не гарантирует, что первая рекомендация является лучшим агрономическим выбором. Корректная формулировка результата --- explainable decision-support prototype для next-crop suggestion.

\subsection{Почему нужен рекомендательный формат}

Обычная классификационная постановка сводит результат к одному классу, но для прикладной задачи выбора культуры такой формат слишком жесткий. Даже если Top-1 прогноз ошибается, фактическая культура может находиться среди нескольких первых кандидатов. В этом случае модель все еще полезна: она сужает пространство вариантов и дает пользователю shortlist для анализа.

Именно поэтому центральным результатом работы является Top-k workflow. Если система показывает только один вариант, она ведет себя как классификатор. Если она показывает несколько вероятных культур, confidence bucket и объяснительный контекст, она становится инструментом поддержки принятия решений. Такой формат честнее отражает неопределенность данных и лучше подходит для областей, где окончательное решение должно оставаться за специалистом.

Рекомендательная постановка также делает более осмысленным анализ трудных классов. Для некоторых культур Top-1 recall ниже, чем для массовых классов, однако Top-3 и Top-5 покрытие существенно выше. Поэтому shortlist особенно важен там, где одна жесткая рекомендация была бы недостаточно надежной.

\subsection{Позиционирование относительно существующих решений}

Данная работа не ставит целью конкурировать с широкими farm management или agronomy platforms. Такие системы обычно охватывают учет операций, экономику хозяйства, планирование ресурсов, погодные и почвенные данные, интеграции с оборудованием и другие производственные процессы. В дипломе решается более узкая и проверяемая задача: построить воспроизводимый прототип, который по историческим данным о культурах формирует вероятностные рекомендации следующей культуры.

Такое сужение задачи является принципиальным. Оно позволяет явно зафиксировать источник данных, split contract, class mapping, модельные артефакты и метрики. В результате работа оценивается не по широте продуктовой функциональности, а по качеству исследовательского контура: подготовка данных, обучение финальной модели, Top-k оценка, confidence analysis, сравнение с Markov baseline и объяснимость.

\subsection{Границы и ограничения постановки}

Границы постановки важны для корректной интерпретации результатов. Во-первых, модель обучается на исторических закономерностях и не доказывает причинную агрономическую оптимальность культуры. Во-вторых, в данных не учитываются экономические факторы, прогноз погоды, свойства почвы и управленческие решения хозяйства. В-третьих, редкие классы исключены из финальной постановки по порогу 1\%, что улучшает устойчивость модели, но ограничивает покрытие менее распространенных культур.

Кроме того, knowledge graph используется как semantic explainability layer, а не как замена основной модели. Spatial demo показывает, как рекомендации могут быть связаны с картой, но не является отдельной предиктивной системой. Эти ограничения не обесценивают работу, а задают честную рамку: прототип помогает анализировать вероятные варианты, но не заменяет агрономическую экспертизу.

\section{Архитектура решения}

\subsection{Общая архитектура системы}

Архитектура решения состоит из нескольких слоев. Первый слой отвечает за данные: исходные CSB/CDL таблицы, идентификаторы полей, годовые классы культур, площадь и координаты. Второй слой выполняет предобработку: сопоставление CDL-кодов, укрупнение культур, фильтрацию нерелевантных классов и построение оконной выборки. Третий слой содержит финальное предиктивное ядро --- \texttt{baseline} CatBoost.

Поверх модели построен recommendation layer. Он преобразует вероятности классов в Top-k список кандидатов и считает recommendation-метрики. Далее добавляется confidence layer, который разделяет случаи по уровню надежности. После этого knowledge graph и rule layer дают объяснительный контекст для кандидатов: переходную статистику, региональную типичность и support/warning clues. Завершающий слой --- spatial demo, показывающий рекомендации и объяснения в привязке к полям на карте.

\subsection{Notebook-first pipeline}

Репозиторий организован как notebook-first pipeline. Это соответствует исследовательскому характеру диплома: каждый этап выполняется в отдельном ноутбуке, сохраняет durable artifacts и передает их следующему этапу. Такой workflow удобен для воспроизводимости, потому что тяжелые промежуточные результаты не требуется пересчитывать без необходимости.

Основная цепочка включает следующие этапы:
\begin{enumerate}
    \item \texttt{01\_data\_analysis\_and\_prep.ipynb} --- загрузка, очистка, сопоставление CDL-кодов и базовая фильтрация;
    \item \texttt{02\_window\_dataset\_and\_feature\_table.ipynb} --- построение окон \texttt{history -> target} и фильтрация редких target-классов;
    \item \texttt{03\_baselines\_and\_split.ipynb} --- фиксированный group-aware split по \texttt{CSBID}, class mapping и простые baseline-ориентиры;
    \item \texttt{04\_catboost\_baseline.ipynb} --- обучение и сохранение финальной CatBoost baseline-модели;
    \item \texttt{05\_catboost\_improved.ipynb} --- экспериментальная модель с derived-признаками;
    \item \texttt{06\_catboost\_plus\_graph.ipynb} --- проверка гибрида CatBoost и transition graph;
    \item \texttt{07\_recommendation\_confidence.ipynb} --- Top-k рекомендации, confidence buckets и калибровка;
    \item \texttt{08\_markov\_recommendation\_eval.ipynb} --- Markov-1 baseline и сравнение с CatBoost;
    \item \texttt{09\_knowledge\_graph\_explainability.ipynb} --- knowledge graph и rule-based объяснимость;
    \item \texttt{10\_spatial\_explainable\_map.ipynb} --- пространственная демонстрационная карта;
    \item \texttt{11\_spatial\_bbox\_demo\_selection.ipynb} --- локальная bbox-выборка для компактной интерактивной демонстрации.
\end{enumerate}

Повторяемая логика вынесена в \texttt{notebooks/\_shared\_notebook\_utils.py} и \texttt{notebooks/\_shared\_geo\_demo\_utils.py}. Это снижает риск расхождения между этапами и сохраняет notebook-first workflow без превращения проекта в production package.

\subsection{Ключевые проектные контракты}

Проект является artifact-driven: активная модель, split, mapping классов и feature contract фиксируются в артефактах и metadata-файлах. Для активной модели источником истины является \texttt{artifacts/model\_registry.json}. В текущем состоянии активная и финальная модель --- \texttt{baseline}. Метаданные финальной модели лежат в \texttt{artifacts/catboost/baseline/meta.json}; split и class mapping зафиксированы в \texttt{artifacts/research\_checkpoint/baseline\_meta.json}.

Критически важным является group-aware split по \texttt{CSBID}. Записи одного поля не должны попадать одновременно в разные части выборки, иначе оценка качества может быть завышена за счет утечки между train и test. Также фиксируется \texttt{random\_state=42}. Изменение split, class mapping, порога редких классов, feature contract или активной registry-модели является breaking change и в данном черновике не выполняется.

\section{Данные и подготовка выборки}

\subsection{Источник данных}

В работе используются данные USDA NASS CSB/CDL за 2017--2024 годы. CDL предоставляет годовые классы культур, а CSB используется как источник полевых идентификаторов, площади и пространственного контекста. Каждая запись содержит последовательность культур по годам и дополнительные признаки поля, включая \texttt{CSBID}, \texttt{CSBACRES}, \texttt{STATEFIPS}, \texttt{CNTYFIPS}, \texttt{INSIDE\_X} и \texttt{INSIDE\_Y}.

На первом этапе данные приводятся к виду, пригодному для моделирования. CDL-коды сопоставляются с понятными названиями и укрупненными группами культур. Нерелевантные для целевой постановки классы удаляются. После базовой очистки сохранен артефакт \texttt{artifacts/data/01\_df\_prepared\_filtered.pkl}. По \texttt{artifacts/data/01\_meta.json} он содержит 8 110 870 строк, 33 столбца и 8 годовых CDL-столбцов.

\subsection{Предобработка данных}

Предобработка решает две задачи. Первая --- сделать исходные классы интерпретируемыми для дальнейшего анализа. Для этого коды CDL переводятся в названия культур и укрупненные группы. Вторая --- отделить классы, не соответствующие целевой agricultural rotation постановке. Это необходимо, чтобы финальная модель работала с осмысленным набором культур, а не с техническими или нерелевантными категориями.

Результатом является очищенная таблица, которая сохраняет связь с полями и годами, но уже имеет признаки, удобные для построения оконной выборки. В дальнейшем именно этот артефакт используется как вход для второго этапа pipeline.

\subsection{Формирование оконной выборки}

Из годовых последовательностей формируется оконный датасет. Окно имеет вид:
\begin{center}
\texttt{history\_1, history\_2, history\_3 -> target}
\end{center}
Три последовательных года используются как история, а следующий год становится целевой культурой. Поскольку доступны годы с 2017 по 2024, из одной исходной записи можно получить несколько перекрывающихся окон.

Такой формат соответствует постановке next-crop recommendation. Модель получает не только последнюю культуру, но и более длинный контекст. Это важно, потому что повторяемость, чередование и региональные особенности могут проявляться не в одном году, а в последовательности.

Полный оконный датасет содержит 40 554 350 строк. После дальнейшей фильтрации редких target-классов остается 38 511 210 строк и 11 столбцов.

\subsection{Фильтрация редких классов}

После построения окон выполняется фильтрация редких целевых классов. Порог редкости равен 1\% от общего числа наблюдений. Это решение направлено на устойчивость модели и корректность оценки: слишком малые классы усиливают дисбаланс, усложняют обучение и делают метрики менее надежными.

После фильтрации остается 9 целевых классов: \texttt{corn}, \texttt{cotton}, \texttt{fallow}, \texttt{forage\_hay}, \texttt{legumes}, \texttt{other\_cereals}, \texttt{sorghum}, \texttt{soybeans}, \texttt{wheat}. Удаленные редкие классы зафиксированы в \texttt{artifacts/data/02\_meta.json}: \texttt{oilseeds\_other}, \texttt{sugar\_crops}, \texttt{peanuts}, \texttt{rice}, \texttt{root\_tubers}, \texttt{sunflower}, \texttt{specialty\_crops}, \texttt{vegetables\_melons}, \texttt{fruits\_berries}.

\subsection{Разделение на train, validation и test}

Разделение на train, validation и test выполнено один раз и далее переиспользуется во всех этапах. Split является group-aware по \texttt{CSBID}: записи одного поля не должны попадать в разные части выборки. Это защищает оценку от утечки и делает сравнение моделей корректным.

Используется \texttt{random\_state=42}. Фактические размеры выборок:
\begin{itemize}
    \ITEM train: 26 957 630 строк;
    \ITEM validation: 5 776 566 строк;
    \ITEM test: 5 777 014 строк.
\end{itemize}
Class mapping содержит 9 классов и зафиксирован в \texttt{artifacts/research\_checkpoint/baseline\_meta.json}.

\subsection{Признаки финальной модели}

Финальная baseline-модель CatBoost использует компактный набор признаков. Категориальные признаки: \texttt{history\_1}, \texttt{history\_2}, \texttt{history\_3}, \texttt{STATEFIPS}, \texttt{CNTYFIPS}. Числовые признаки: \texttt{CSBACRES}, \texttt{INSIDE\_X}, \texttt{INSIDE\_Y}.

Такой feature contract отражает основную гипотезу работы: следующая культура зависит не только от последней культуры, но и от трехлетней истории, регионального контекста и характеристик поля. Одновременно набор признаков остается достаточно небольшим и воспроизводимым, что важно для дипломного исследовательского прототипа.

\section{Базовые модели и финальная CatBoost-модель}

\subsection{Простые baseline-подходы}

Простые baseline-подходы нужны для проверки того, дает ли более сложная модель реальный прирост качества. В проекте используются ориентиры вроде наиболее частого класса, правила последней культуры и Markov-1 модели. Они не предназначены для финального применения, но задают нижнюю границу качества и помогают интерпретировать результат CatBoost.

Наиболее важным для рекомендательной оценки является Markov-1 baseline. Он использует только последнюю культуру и оценивает переходы вида:
\begin{center}
\texttt{P(target | history\_3)}.
\end{center}
Такой baseline интерпретируем, но ограничен: он не учитывает более длинную историю, регион, площадь и координаты.

\subsection{CatBoost baseline}

Финальной моделью диплома является \texttt{baseline} CatBoost. CatBoost выбран как модель, хорошо работающая с табличными данными и категориальными признаками. В данной задаче категориальными являются история культур и региональные коды, поэтому использование CatBoost позволяет работать с ними без чрезмерной ручной обработки.

На выходе четвертого этапа сохраняются \texttt{artifacts/catboost/baseline/model.cbm} и \texttt{artifacts/catboost/baseline/meta.json}. Метаданные фиксируют feature contract, categorical features, numeric features и class mapping. Это позволяет downstream-этапам использовать модель без неявных предположений.

\subsection{Экспериментальная ветка improved}

В отдельной ветке была проверена \texttt{improved}-модель с derived-признаками истории: количеством уникальных культур, наличием \texttt{fallow}, наличием \texttt{legumes}, повторяемостью соседних лет и повтором всей истории. Эта ветка полезна как исследовательский эксперимент, показывающий возможное направление развития признакового пространства.

Однако \texttt{improved} не является финальной моделью диплома. Проектный контракт и registry фиксируют \texttt{baseline} как активную модель. Поэтому в выводах и основном тексте \texttt{improved} описывается только как experimental non-final branch.

\subsection{Результаты классификационной оценки}

Финальная baseline-модель на test показывает Accuracy@1 = 0.69901, Macro F1 = 0.56849 и Weighted F1 = 0.68963. На validation значения близки: Accuracy@1 = 0.69928, Macro F1 = 0.56877 и Weighted F1 = 0.68997. Близость validation и test указывает на стабильность оценки на зафиксированном split.

\begin{center}
\begin{longtable}{|p{4cm}|p{3cm}|p{3cm}|p{3cm}|}
\caption{Классификационные метрики CatBoost baseline}\\
\hline
Выборка & Accuracy@1 & Macro F1 & Weighted F1\\
\hline
Validation & 0.69928 & 0.56877 & 0.68997\\
\hline
Test & 0.69901 & 0.56849 & 0.68963\\
\hline
\end{longtable}
\end{center}

Macro F1 ниже, чем Weighted F1, потому что одинаково учитывает все классы, включая трудные и менее представленные. Это ожидаемо для несбалансированной agricultural classification задачи и дополнительно показывает, почему важно анализировать не только агрегированную Top-1 метрику.

\subsection{Гибрид CatBoost и transition graph}

В отдельном эксперименте была проверена идея добавить к вероятностям CatBoost графовый сигнал переходов культур:
\begin{equation}
    score = \alpha \cdot p_{\mathrm{catboost}} + \beta \cdot p_{\mathrm{graph}}.
\end{equation}
Проверялись веса $(0.9, 0.1)$, $(0.8, 0.2)$ и $(0.7, 0.3)$. Лучший вариант на test имел $\alpha=0.9$ и $\beta=0.1$, но показал Accuracy = 0.69871, Macro F1 = 0.56439 и Weighted F1 = 0.68873, то есть немного уступил чистому CatBoost.

Этот результат важен методически. Он показывает, что простой graph reranking не следует делать финальной моделью, но графовый сигнал можно использовать как объяснительный слой для анализа кандидатов.

\section{Recommendation layer и оценка качества}

\subsection{Переход от классификации к рекомендации}

После обучения CatBoost вероятности классов преобразуются в ранжированный список кандидатов. Это меняет прикладной смысл системы: она перестает быть только классификатором и становится next-crop recommendation workflow. Пользователь получает не одно предписание, а несколько вероятных вариантов.

В дипломе используются метрики Accuracy@1, Hit@3 и Hit@5. Accuracy@1 совпадает с обычной точностью top-1 прогноза. Hit@3 показывает, попала ли фактическая культура в три первых кандидата, а Hit@5 --- в пять первых кандидатов.

\subsection{Top-k результаты CatBoost}

Для baseline CatBoost на test получены Accuracy@1 = 0.69901, Hit@3 = 0.95124 и Hit@5 = 0.99082. Это ключевой результат диплома: при выдаче одного варианта правильная культура совпадает примерно в 69.9\% случаев, но при выдаче трех кандидатов она находится в списке примерно в 95.1\% случаев, а при пяти кандидатах --- примерно в 99.1\% случаев.

\begin{center}
\begin{longtable}{|p{4cm}|p{3cm}|p{3cm}|p{3cm}|}
\caption{Top-k метрики CatBoost baseline}\\
\hline
Выборка & Accuracy@1 & Hit@3 & Hit@5\\
\hline
Validation & 0.69928 & 0.95119 & 0.99079\\
\hline
Test & 0.69901 & 0.95124 & 0.99082\\
\hline
\end{longtable}
\end{center}

Такой разрыв между Top-1 и Top-k качеством подтверждает, что система особенно сильна как shortlist-рекомендатель. Это также поддерживает центральное позиционирование работы как decision-support prototype.

\subsection{Сравнение CatBoost и Markov-1}

Для проверки пользы финальной модели CatBoost сравнивается с Markov-1 baseline. На test Markov-1 показывает Accuracy@1 = 0.55722, Hit@3 = 0.85530 и Hit@5 = 0.93172. CatBoost превосходит его по всем трем метрикам.

\begin{center}
\begin{longtable}{|p{5cm}|p{3cm}|p{3cm}|p{3cm}|}
\caption{Сравнение CatBoost и Markov-1 на test}\\
\hline
Модель & Accuracy@1 & Hit@3 & Hit@5\\
\hline
CatBoost baseline & 0.69901 & 0.95124 & 0.99082\\
\hline
Markov-1 baseline & 0.55722 & 0.85530 & 0.93172\\
\hline
\end{longtable}
\end{center}

Сравнение показывает, что трехлетняя история и контекстные признаки действительно добавляют полезную информацию. Модель не просто воспроизводит частые переходы после последней культуры, а использует более широкий контекст поля и региона.

\subsection{Поклассовый анализ}

Поклассовые результаты показывают неоднородность качества. На test наиболее устойчивые классы по recall@1: \texttt{forage\_hay} --- 0.85049, \texttt{cotton} --- 0.76594, \texttt{soybeans} --- 0.76050, \texttt{wheat} --- 0.73038, \texttt{corn} --- 0.71866. Более трудные классы: \texttt{fallow} --- 0.40681, \texttt{sorghum} --- 0.30611, \texttt{other\_cereals} --- 0.24006, \texttt{legumes} --- 0.13845.

При этом Top-k формат частично компенсирует сложность таких классов. Например, для \texttt{legumes} recall@1 составляет 0.13845, но recall@3 равен 0.73309, а recall@5 --- 0.90956. Для \texttt{fallow} recall@1 равен 0.40681, но recall@3 достигает 0.80443. Это демонстрирует, почему система должна использоваться как рекомендательная, а не как жесткая замена экспертному выбору.

\section{Уверенность рекомендаций и калибровка}

\subsection{Confidence buckets}

В рекомендательной системе важно понимать не только состав shortlist, но и надежность рекомендации. Поэтому в проекте введены три confidence buckets: \texttt{high\_confidence}, \texttt{medium\_confidence} и \texttt{low\_confidence}. Они позволяют разделять случаи, где модель выражает сильную уверенность, и случаи, где пользователю следует особенно внимательно рассмотреть альтернативы.

Этот слой не отменяет рекомендации. Даже при низкой уверенности модель может показывать полезный shortlist, но интерфейс и текст интерпретации должны явно сообщать, что случай является спорным.

\subsection{Результаты по confidence buckets}

На test распределение по confidence buckets выглядит следующим образом. В группе high confidence содержится 3 042 231 наблюдение; Accuracy@1 = 0.84147, Hit@3 = 0.97713. В группе medium confidence содержится 1 691 260 наблюдений; Accuracy@1 = 0.60931, Hit@3 = 0.94252. В группе low confidence содержится 1 043 523 наблюдения; Accuracy@1 = 0.42902, Hit@3 = 0.88987.

\begin{center}
\begin{longtable}{|p{4cm}|p{3cm}|p{3cm}|p{3cm}|}
\caption{Качество по confidence buckets на test}\\
\hline
Группа & Наблюдения & Accuracy@1 & Hit@3\\
\hline
High confidence & 3 042 231 & 0.84147 & 0.97713\\
\hline
Medium confidence & 1 691 260 & 0.60931 & 0.94252\\
\hline
Low confidence & 1 043 523 & 0.42902 & 0.88987\\
\hline
\end{longtable}
\end{center}

Результаты показывают, что confidence layer действительно разделяет случаи по надежности. В high confidence модель значительно точнее среднего уровня, а в low confidence Top-1 качество ниже. При этом даже в low confidence Hit@3 остается высоким, что поддерживает идею показывать несколько кандидатов вместо одного.

\subsection{Калибровка вероятностей}

Дополнительно проверялась калибровка вероятностей на preview-выборках по 5000 объектов. На validation ECE = 0.02305, Brier = 0.16440, mean confidence = 0.75026, Accuracy@1 = 0.73560. На test ECE = 0.01737, Brier = 0.16120, mean confidence = 0.75398, Accuracy@1 = 0.74800.

Низкий ECE на preview-выборке означает, что вероятности в целом согласованы с фактической точностью. В рамках данного прототипа это поддерживает использование вероятностей не только для сортировки кандидатов, но и как понятного confidence-сигнала.

\section{Explainability через knowledge graph}

\subsection{Зачем нужен explainability layer}

Top-k список сам по себе не объясняет, почему кандидаты оказались в рекомендации. Для поддержки принятия решений полезно показать, какие факторы поддерживают или ослабляют конкретный вариант. Поэтому после предиктивной модели построен knowledge graph explainability layer.

Этот слой не заменяет CatBoost и не доказывает агрономическую правильность рекомендации. Его задача --- дать семантический и rule-based контекст: переходы между культурами, принадлежность к группам, региональную типичность и предупреждения по простым правилам.

\subsection{Структура knowledge graph}

По \texttt{artifacts/results/knowledge\_graph\_explainability/run\_meta.json} граф содержит 2 827 узлов, 18 154 ребра, 5 правил, 4 case studies и 4 графические визуализации кейсов. В графе представлены культуры, группы культур, регионы и правила. Ребра отражают переходы между культурами по train-only статистике, принадлежность культуры к укрупненной группе, типичность культуры для региона и связи правил с кандидатами.

Такой граф следует понимать как табличный property-graph для объяснения, а не как внешнюю экспертную базу знаний. Он строится поверх текущих артефактов и помогает интерпретировать shortlist.

\subsection{Rule layer}

В explainability layer включены пять правил:
\begin{itemize}
    \ITEM \texttt{repeat\_crop\_warning} --- предупреждение о повторе культуры;
    \ITEM \texttt{same\_group\_saturation} --- предупреждение о насыщении одной группой культур;
    \ITEM \texttt{legume\_break\_bonus} --- поддержка бобовых как break crop;
    \ITEM \texttt{fallow\_break\_effect} --- интерпретация эффекта fallow;
    \ITEM \texttt{region\_typicality} --- поддержка регионально типичных культур.
\end{itemize}

Правила дают support/warning clues и не являются полноценной экспертной агрономической системой. Их назначение --- сделать объяснение более читаемым и показать, почему отдельный кандидат может быть согласован или не согласован с простым графовым контекстом.

\subsection{Candidate-level output}

Для каждого кандидата может быть собран объяснительный набор признаков: вероятность модели, позиция в shortlist, графовый support, уровень поддержки, предупреждения и альтернативы. Важным полем является \texttt{graph\_support\_score}. Он не является вероятностью и не должен интерпретироваться как доказательство причинной пользы культуры. Корректная интерпретация: насколько кандидат согласован с переходной статистикой, региональной типичностью и rule-based контекстом.

\subsection{Интерпретация согласованности}

В case studies распределение согласованности top-1 с графом получилось следующим: один кейс имеет \texttt{strong} alignment, три кейса --- \texttt{weak} alignment. Категория \texttt{weak} не означает, что модель ошиблась. Она означает, что graph/rule evidence слабее поддерживает top-1 кандидата или показывает сильные альтернативы внутри shortlist.

Это важный результат для объяснимости. Граф не обязан всегда подтверждать CatBoost. Иногда он помогает обнаружить спорные случаи и сделать рекомендацию более прозрачной для пользователя.

\section{Case studies и пространственная демонстрация}

\subsection{Case studies}

В проекте подготовлены четыре case studies: high confidence, medium confidence, low confidence и difficult class. Каждый кейс предназначен для анализа рекомендации на уровне отдельного объекта: фактический target, top-1 выбор модели, вероятность модели, графовая поддержка и согласованность top-1 с graph/rule контекстом.

\begin{center}
\begin{longtable}{|p{3.8cm}|p{2.2cm}|p{2.2cm}|p{2.2cm}|p{2.8cm}|}
\caption{Краткое описание case studies}\\
\hline
Кейс & Target & Top-1 & Score top-1 & Alignment\\
\hline
High confidence & \texttt{forage\_hay} & \texttt{forage\_hay} & 0.92991 & strong\\
\hline
Medium confidence & \texttt{corn} & \texttt{corn} & 0.66878 & weak\\
\hline
Low confidence & \texttt{legumes} & \texttt{legumes} & 0.43222 & weak\\
\hline
Difficult class & \texttt{soybeans} & \texttt{corn} & 0.68867 & weak\\
\hline
\end{longtable}
\end{center}

В high confidence кейсе top-1 модели совпадает с фактическим target \texttt{forage\_hay}; графовая поддержка top-1 также является самой сильной среди кандидатов. Это пример согласованного простого случая, где модельное ранжирование и graph evidence поддерживают один и тот же вариант.

В medium confidence кейсе top-1 \texttt{corn} также совпадает с фактическим target, но graph/rule слой сильнее поддерживает альтернативу \texttt{soybeans}. Это не делает прогноз ошибочным, но показывает, что внутри shortlist есть семантически значимая альтернатива. В low confidence кейсе target \texttt{legumes} также находится на первом месте модели, однако графовая поддержка для него слабая, а \texttt{soybeans} поддерживается заметно сильнее. Такой кейс показывает, зачем нужен confidence и explainability layer: пользователь должен видеть не только ответ модели, но и степень согласованности этого ответа с дополнительным контекстом.

В difficult class кейсе top-1 модели равен \texttt{corn}, а фактический target --- \texttt{soybeans}. При этом \texttt{soybeans} присутствует в shortlist на третьей позиции и имеет более сильную graph support оценку, чем top-1. Этот пример особенно важен для защиты Top-k постановки: при оценке только top-1 случай считался бы ошибкой, но в рекомендательном режиме правильный кандидат остается доступен пользователю вместе с объяснительным сигналом.

\subsection{Spatial explainable demo}

На поздних этапах добавлена пространственная демонстрация. Она связывает рекомендации с геометрией полей и показывает, как результат может выглядеть в пользовательском интерфейсе. Пользователь видит поле на карте, открывает подсказку и получает Top-k кандидатов, confidence и объяснительный контекст.

Важно, что spatial demo не является новой моделью и не меняет метрики. Это демонстрационный и аналитический слой поверх уже построенного recommendation workflow. Финальным предиктивным ядром остается \texttt{baseline} CatBoost.

\subsection{BBox demo selection}

Для повышения читаемости карты добавлена локальная bbox-выборка. По \texttt{spatial\_layer\_comparison\_summary.csv} полный региональный слой содержит 19 919 полей, слой после фильтрации 9 целевых классов содержит 16 663 поля, то есть сохраняет 83.7\% полей полного регионального слоя. Демонстрационный subset содержит 20 полей.

По \texttt{spatial\_bbox\_selection\_summary.csv} локальный bbox-контекст содержит 157 полей, с test-строками сопоставлено 23 поля, и итоговый clickable demo dataset содержит 23 поля. Это делает карту компактной и пригодной для демонстрации без пересчета модели и без изменения основных результатов.

\section{Ограничения и направления развития}

\subsection{Ограничения}

Основные ограничения текущей версии следующие:
\begin{itemize}
    \ITEM модель обучается на исторических закономерностях и не доказывает агрономическую оптимальность культуры;
    \ITEM в данных не учитываются экономические факторы, погодные прогнозы, почвенные ограничения и управленческие решения хозяйства;
    \ITEM редкие классы исключены из финальной постановки по порогу 1\%;
    \ITEM на трудных классах Top-1 качество заметно ниже, чем на массовых;
    \ITEM гибрид CatBoost и transition graph не улучшил baseline-метрики;
    \ITEM KG explainability является rule-based интерпретацией, а не причинной моделью;
    \ITEM spatial map является демонстрационным интерфейсом, а не отдельной системой принятия решений.
\end{itemize}

Эти ограничения важны для корректного позиционирования работы. Система должна применяться как аналитический и рекомендательный слой, а не как автоматический источник окончательного агрономического решения.

\subsection{Направления развития}

Дальнейшее развитие может идти в нескольких направлениях. Во-первых, knowledge graph можно расширить за счет expert-in-the-loop редактирования правил и связей. Во-вторых, модель можно дополнить внешними признаками, если они будут надежно доступны и согласованы с постановкой: почвенными характеристиками, погодными данными, экономическими индикаторами или управленческими ограничениями. В-третьих, spatial interface можно развить из демонстрационной карты в более удобный аналитический интерфейс для просмотра полей, рекомендаций и объяснений.

При этом любое расширение должно сохранять воспроизводимость и не нарушать текущие проектные контракты без явного пересмотра: split по \texttt{CSBID}, class mapping, feature contract и активную модель registry.

\specialsection{Заключение}

В работе построен первый полный контур объяснимого исследовательского прототипа рекомендательной системы следующей культуры по истории поля. Данные USDA NASS CSB/CDL за 2017--2024 годы были подготовлены, преобразованы в оконную выборку, очищены от редких целевых классов и разделены на train, validation и test с group-aware split по \texttt{CSBID}.

Финальной моделью диплома является \texttt{baseline} CatBoost. Она использует трехлетнюю историю культур, региональные признаки, площадь поля и координаты. На test модель показала Accuracy@1 = 0.69901, Macro F1 = 0.56849 и Weighted F1 = 0.68963. В рекомендательном режиме получены Hit@3 = 0.95124 и Hit@5 = 0.99082, что подтверждает ценность Top-k shortlist по сравнению с одиночным прогнозом.

Сравнение с Markov-1 baseline показало, что CatBoost существенно превосходит простую переходную модель по Accuracy@1, Hit@3 и Hit@5. Confidence analysis разделил рекомендации по уровню надежности, а knowledge graph добавил семантический слой объяснимости через переходы, региональную типичность и rule-based support/warning clues. Пространственная демонстрация показала, как рекомендации могут быть связаны с картой полей.

Итоговую систему корректно описывать как explainable Top-k decision-support prototype for next-crop suggestion. Она предлагает вероятные культуры, оценивает уверенность и объясняет контекст рекомендации, но не заменяет экспертную агрономическую проверку и не является гарантированным оптимизатором севооборота.

\end{document}
